%=====================================================

% lista de símbolos

\begin{listasimb}

\begin{longtable}[l]{p{0.5\linewidth}p{15.0\linewidth}}
$\mathbb{C}$ - Matriz constitutiva elástica \\
$\mathbb{D}$ - Matriz constitutiva viscosa \\
$E$ - Módulo de elasticidade \\
$\epsilon$ - Vetor de deformações \\
$\varepsilon$ - Tensor de deformações \\
$t$ - Tempo \\
$b$ - Forças volumétricas \\
$\varphi$ - funções de base \\
$\mathbf{V}$ - Volume  \\
$\mathbf{S}$ - Superfície  \\
$\hat{n}$ - Vetor unitário normal \\
$\mathbf{x}$ - Coordenadas cartesianas em notação compacta \\
$u(\mathbf{x},t)$, $u(t)$ ou $u$ - Deslocamentos em função do espaço e do tempo \\
$\dot{u}$ - Derivada de primeira ordem dos deslocamentos em relação ao tempo ou velocidade \\
$\ddot{u}$ - Derivada de segunda ordem dos deslocamentos em relação ao tempo ou aceleração \\
$u_h$ - Função de deslocamentos aproximada \\
$m$ - massa \\
$\nu$ - Coeficiente de Poisson.  \\
$\tau$ - Tensões de cisalhamento \\ 
$F_{i}$ ou $\{ F \}$ - Vetor de forças generalizado \\
$f$ - Vetor de forçar generalizado elementar \\
$q_n(t)$ - Coordenada generalizada \\
$\phi_n$ - Vetor modo de vibração natural \\
$\omega$ - frequência \\
$\omega_n$ - frequência natural  \\
$\rho$ - Densidade \\
$\mathsf{T}$ - tensor de tensões \\
$R$ - Função resíduo \\
$\mathsf{t}^{(\cdot)}$ - vetor de tensões em relação à direção $(\cdot)$  \\
$\left[ I \right]$ - matriz identidade \\
$\left[ \Omega^2 \right]$ - matriz espectral  \\
$\left[ K \right]$ - matriz de rigidez \\
$\left[ K \right]_E$ - matriz de rigidez local\\
$\left[ K \right]_E^{const}$ - Parcela constitutiva da matriz de rigidez elementar \\
$\left[ K  \right]_E^{stab}$ - Parcela de estabilização da matriz de rigidez global \\
$\left[ M \right]$ - matriz de massa  \\
$\left[ M \right]_E$ - matriz de massa local \\
$\left[ M \right]_E^{const}$ - Parcela constitutiva da matriz de massa elementar \\
$\left[ M \right]_E^{stab}$ - Parcela de estabilização da matriz de massa elementar \\
$\left[ C \right]$ - Matriz de amortecimento \\
$\left[ C \right]_E$ - matriz de amortecimento local \\
$w$ - funções peso\\
$\mathcal{G}^{\nabla}$ - Matriz de Gram $H^1$ para o operador $\Pi_k^{\nabla}$ \\
$\mathcal{G}^{\epsilon}$ - Matriz de Gram $H^1$ para o operador $\Pi_k^{\epsilon}$ \\
$\tilde{G}^\nabla$ - Matriz de Gram $H^1$ com a primeira linha nula para o operador $\Pi_k^{\nabla}$ \\
$\mathcal{B}^\nabla$ - Matriz dos termos virtuais para o operador $\Pi_k^{\nabla}$ \\
$\mathcal{B}^\epsilon$ - Matriz dos termos virtuais para o operador $\Pi_k^{\epsilon}$ \\
$\mathcal{D}^\nabla$ - Matriz de transformação de bases para o operador $\Pi_k^{\nabla}$\\
$\mathcal{H}$ - Matriz de Gram $L^2$\\
$\mathcal{C}$ - Matriz dos termos virtuais para o operador $\Pi_k^0$\\
$e_{abs}$ - Erro absoluto \\
$e_{NRMSE}$ - Raiz normada do erro médio quadrático\\
$e_{MG}$ - Erro da média geométrica \\
$|E|$ - Área de um elemento
$h_E$ - Diâmetro de um elemento\\
$\Omega$ - Domínio\\
$\Omega_v$ - Domínio de um elemento\\
$V$ - vértice\\
$\mathbb{V}(\Omega)$ - Espaço de solução no domínio $\Omega$\\
$\mathbb{V}_k(\Omega)$ - Espaço de solução aproximado por um polinômio de grau $k$\\  
$\mathbb{P}_k$ - Espaço de funções polinomiais de grau $k$\\
$\mathbb{B}$ - Espaço de funções aproximadas auxiliar associado ao contorno do elemento\\
$p_\alpha$ - conjunto de funções polinomiais\\
$n_v$ - número de vértices\\
$N_k^P$ - Vetor de expansão polinomial de grau $k$\\
$H_k^P$ - matriz de expansão polinomial de grau $k$\\
$n_e$ - número de arestas\\
$n_{dof}$ - número de graus de liberdade \\
$v$ - função teste\\
$v_h$ função teste aproximada\\ 
$\varphi$ - função de base \\
$k$ - grau polinomial \\
$\Pi_{k}^\nabla$ - Projetor $H^1$ de grau $k$\\
$\Pi_{k,*}^0$ - Projetor $L^2$ de grau $k$ na base polinomial\\
$\Pi_{k}^0$ - Projetor $L^2$ de grau $k$ na base canônica\\
$\Pi_{k-1}^\epsilon$ - Projetor $H^1$ de grau $k-1$\\
$B^{(\nabla)}$ - matriz com derivadas do operador gradiente \\
$B^{(\epsilon)}$ - matriz com derivadas do operador de deformações \\
$n_k$ - dimensão de um polinômio de grau $k$\\
$\xi$ e $\eta$ - Coordenas locais\\
$e$ - aresta\\
$\{a\}$ - Vetor de coeficientes\\
$P_0$ - operador sobre constantes\\
$N_E$ - matriz de componentes normais\\


\end{longtable}

\end{listasimb}

%=====================================================
